\textbf{Duración:} 3 horas (en aula)\\
\textbf{Enfoque:} De objetivos de aprendizaje a arco narrativo y storyboard de una unidad o conferencia técnica.

\subsubsection*{Objetivos de la sesión}
Al finalizar, la persona participante:
\begin{enumerate}
\item Convierte objetivos de una unidad en mensajes clave y arco narrativo.
\item Produce un storyboard de 8--12 nodos con gancho, desarrollo y cierre.
\item Anticipa malentendidos típicos del tema STEM.
\item Justifica decisiones de carga cognitiva (qué va en slide vs.\ en notas).
\end{enumerate}

\subsubsection*{Agenda (resumen)}
\begin{tabularx}{\textwidth}{@{}lclX@{}}
\toprule
Bloque & Tiempo & Contenido \\
\midrule
Recap & 10 min & De texto largo a experiencia de clase \\
Narrativa & 30 min & Objetivos $\rightarrow$ mensajes $\rightarrow$ arco \\
Demo storyboard & 20 min & 8--12 nodos \\
Taller & 75 min & Storyboard de unidad EIC \\
Pares & 25 min & Checklist narrativa \\
Cierre & 10 min & TI hacia deck \\
\bottomrule
\end{tabularx}

\subsubsection*{Producto esperado}
Storyboard narrativo + justificación de mensajes clave (E5).

\subsubsection*{Trabajo independiente relacionado}
Sesión 6: deck MV, notas de orador y accesibilidad. TI: refinar storyboard (bloque C).
